% Generated from spice_optical_depth.md with Pandoc; retain the Markdown source
% as the editable version.
% Options for packages loaded elsewhere
\PassOptionsToPackage{unicode}{hyperref}
\PassOptionsToPackage{hyphens}{url}
\PassOptionsToPackage{dvipsnames,svgnames,x11names}{xcolor}
\documentclass[
]{article}
\usepackage{xcolor}
\usepackage[margin=1in]{geometry}
\usepackage{amsmath,amssymb}
\setcounter{secnumdepth}{-\maxdimen} % remove section numbering
\usepackage{iftex}
\ifPDFTeX
  \usepackage[T1]{fontenc}
  \usepackage[utf8]{inputenc}
  \usepackage{textcomp} % provide euro and other symbols
\else % if luatex or xetex
  \usepackage{unicode-math} % this also loads fontspec
  \defaultfontfeatures{Scale=MatchLowercase}
  \defaultfontfeatures[\rmfamily]{Ligatures=TeX,Scale=1}
\fi
\usepackage{lmodern}
\ifPDFTeX\else
  % xetex/luatex font selection
\fi
% Use upquote if available, for straight quotes in verbatim environments
\IfFileExists{upquote.sty}{\usepackage{upquote}}{}
\IfFileExists{microtype.sty}{% use microtype if available
  \usepackage[]{microtype}
  \UseMicrotypeSet[protrusion]{basicmath} % disable protrusion for tt fonts
}{}
\makeatletter
\@ifundefined{KOMAClassName}{% if non-KOMA class
  \IfFileExists{parskip.sty}{%
    \usepackage{parskip}
  }{% else
    \setlength{\parindent}{0pt}
    \setlength{\parskip}{6pt plus 2pt minus 1pt}}
}{% if KOMA class
  \KOMAoptions{parskip=half}}
\makeatother
\usepackage{listings}
\newcommand{\passthrough}[1]{#1}
\lstset{defaultdialect=[5.3]Lua}
\lstset{defaultdialect=[x86masm]Assembler}
\lstset{breaklines=true,breakatwhitespace=false,columns=fullflexible}
\usepackage{longtable,booktabs,array}
\newcounter{none} % for unnumbered tables
\usepackage{calc} % for calculating minipage widths
% Correct order of tables after \paragraph or \subparagraph
\usepackage{etoolbox}
\makeatletter
\patchcmd\longtable{\par}{\if@noskipsec\mbox{}\fi\par}{}{}
\makeatother
% Allow footnotes in longtable head/foot
\IfFileExists{footnotehyper.sty}{\usepackage{footnotehyper}}{\usepackage{footnote}}
\makesavenoteenv{longtable}
\setlength{\emergencystretch}{3em} % prevent overfull lines
\providecommand{\tightlist}{%
  \setlength{\itemsep}{0pt}\setlength{\parskip}{0pt}}
\usepackage{bookmark}
\IfFileExists{xurl.sty}{\usepackage{xurl}}{} % add URL line breaks if available
\urlstyle{same}
\hypersetup{
  pdftitle={SPICE line optical depths from MURaM},
  colorlinks=true,
  linkcolor={Maroon},
  filecolor={Maroon},
  citecolor={Blue},
  urlcolor={Blue},
  pdfcreator={LaTeX via pandoc}}

\title{SPICE line optical depths from MURaM}
\author{}
\date{}

\begin{document}
\maketitle

\passthrough{\lstinline!compute\_spice\_optical\_depth.py!} calculates
the line optical depth through each vertical MURaM column, including
thermal broadening, optional microturbulence, the resolved velocity, and
stimulated emission. It also synthesizes an optically thin comparison
spectrum from the same atomic populations. The reusable physics and
population-table code is in
\passthrough{\lstinline!modules/spiceOpticalDepth.py!}.

This is an \textbf{equilibrium, bound-bound opacity diagnostic}. CHIANTI
populations are not recomputed in response to the radiation trapped by
an optically thick line. Consequently, a large optical depth flags a
failure of the optically thin assumption; the saved comparison intensity
is not a prediction of the emergent intensity in that regime.

\section{Line selection and the instrument
paper}\label{line-selection-and-the-instrument-paper}

The catalogue follows Table 1 of \textbf{SPICE Consortium et al.~(2020),
``The Solar Orbiter SPICE instrument: An extreme UV imaging
spectrometer,'' A\&A 642, A14},
\href{https://doi.org/10.1051/0004-6361/201935574}{DOI},
\href{https://www2.mps.mpg.de/homes/solanki/publikationen/j532.pdf}{accessible
PDF}. The paper specifies bands of 704--790 and 973--1049 Å. Its table
is a selection of diagnostics, not every transition in those bands. Si
XII is a second-order line: the absorbing photon has wavelength 520.67
Å, while its nominal first-order-equivalent detector coordinate is
1041.34 Å. All opacity and emissivity calculations use the physical
photon wavelength.

{\def\LTcaptype{none} % do not increment counter
\begin{longtable}[]{@{}
  >{\raggedright\arraybackslash}p{(\linewidth - 8\tabcolsep) * \real{0.1765}}
  >{\raggedright\arraybackslash}p{(\linewidth - 8\tabcolsep) * \real{0.1765}}
  >{\raggedleft\arraybackslash}p{(\linewidth - 8\tabcolsep) * \real{0.2353}}
  >{\raggedleft\arraybackslash}p{(\linewidth - 8\tabcolsep) * \real{0.2353}}
  >{\raggedright\arraybackslash}p{(\linewidth - 8\tabcolsep) * \real{0.1765}}@{}}
\toprule\noalign{}
\begin{minipage}[b]{\linewidth}\raggedright
CLI identifier
\end{minipage} & \begin{minipage}[b]{\linewidth}\raggedright
Ion
\end{minipage} & \begin{minipage}[b]{\linewidth}\raggedleft
Paper wavelength (Å)
\end{minipage} & \begin{minipage}[b]{\linewidth}\raggedleft
Paper log10 T (K)
\end{minipage} & \begin{minipage}[b]{\linewidth}\raggedright
Default
\end{minipage} \\
\midrule\noalign{}
\endhead
\bottomrule\noalign{}
\endlastfoot
\passthrough{\lstinline!h\_1\_1025!} & H I & 1025.72 & 4.0 & \\
\passthrough{\lstinline!c\_2\_1036!} & C II & 1036.34 & 4.3 & \\
\passthrough{\lstinline!c\_3\_977!} & C III & 977.03 & 4.5 & yes \\
\passthrough{\lstinline!s\_5\_786!} & S V & 786.47 & 5.2 & yes \\
\passthrough{\lstinline!o\_4\_787!} & O IV & 787.72 & 5.2 & yes \\
\passthrough{\lstinline!o\_5\_760!} & O V & 760.43 & 5.4 & yes \\
\passthrough{\lstinline!o\_6\_1031!} & O VI & 1031.93 & 5.5 & yes \\
\passthrough{\lstinline!o\_6\_1037!} & O VI & 1037.64 & 5.5 & yes \\
\passthrough{\lstinline!ne\_6\_1005!} & Ne VI & 1005.79 & 5.6 & \\
\passthrough{\lstinline!si\_7\_1049!} & Si VII & 1049.25 & 5.6 & \\
\passthrough{\lstinline!ne\_8\_770!} & Ne VIII & 770.42 & 5.8 & yes \\
\passthrough{\lstinline!mg\_8\_772!} & Mg VIII & 772.31 & 5.9 & \\
\passthrough{\lstinline!mg\_9\_706!} & Mg IX & 706.02 & 6.0 & yes \\
\passthrough{\lstinline!fe\_10\_1028!} & Fe X & 1028.04 & 6.0 & \\
\passthrough{\lstinline!mg\_11\_997!} & Mg XI & 997.44 & 6.2 & \\
\passthrough{\lstinline!si\_12\_520!} & Si XII, order 2 & 520.67 & 6.3
& \\
\passthrough{\lstinline!fe\_18\_974!} & Fe XVIII & 974.84 & 6.9 & \\
\passthrough{\lstinline!fe\_20\_721!} & Fe XX & 721.55 & 7.0 & \\
\end{longtable}
}

The eight defaults are my practical starting set for this repository's
transition-region/coronal models. They include C III, the O VI doublet,
and Ne VIII, plus intermediate-temperature diagnostics and Mg IX. This
default set is a software choice. The table's approximate temperatures
are labels, \textbf{not temperature masks or substitutes for an
ionization balance}.

Use \passthrough{\lstinline!--all-lines!} for all 18 entries or
\passthrough{\lstinline!--lines!} for a subset. H I is explicitly
labelled an approximation when selected: Ly beta usually needs an NLTE
hydrogen calculation, ideally using the simulation's hydrogen
populations when available. C II and C III also need particular care
with ionization equilibrium and radiative-transfer effects. The script
does not automatically interpret the reader's
\passthrough{\lstinline!H\_1!}, \passthrough{\lstinline!HH!}, or
\passthrough{\lstinline!HHp!} files, whose definitions must first be
established for the particular simulation.

\section{How this fits the existing
repository}\label{how-this-fits-the-existing-repository}

The implementation reuses \passthrough{\lstinline!MuramCube!} from
\passthrough{\lstinline!modules/muram.py!} and follows the transpose
used by \passthrough{\lstinline!MURaMData!}. It opens only the four
required physical cubes. The full-domain input files are:

{\def\LTcaptype{none} % do not increment counter
\begin{longtable}[]{@{}
  >{\raggedright\arraybackslash}p{(\linewidth - 4\tabcolsep) * \real{0.3333}}
  >{\raggedright\arraybackslash}p{(\linewidth - 4\tabcolsep) * \real{0.3333}}
  >{\raggedright\arraybackslash}p{(\linewidth - 4\tabcolsep) * \real{0.3333}}@{}}
\toprule\noalign{}
\begin{minipage}[b]{\linewidth}\raggedright
File
\end{minipage} & \begin{minipage}[b]{\linewidth}\raggedright
Quantity
\end{minipage} & \begin{minipage}[b]{\linewidth}\raggedright
Required unit
\end{minipage} \\
\midrule\noalign{}
\endhead
\bottomrule\noalign{}
\endlastfoot
\passthrough{\lstinline!Header.<iteration>!} & Dimensions, spacings,
simulation time & cm for spacings, s for time \\
\passthrough{\lstinline!eosT.<iteration>!} & Temperature & K \\
\passthrough{\lstinline!eosne.<iteration>!} & Electron density &
cm\^{}-3 \\
\passthrough{\lstinline!result\_prim\_0.<iteration>!} & Mass density & g
cm\^{}-3 \\
\passthrough{\lstinline!result\_prim\_1.<iteration>!} & Native vertical
velocity \passthrough{\lstinline!vx!} & cm s\^{}-1 \\
\end{longtable}
}

An iteration is formatted with at least six digits, for example
\passthrough{\lstinline!400000!}. Iteration increments are not physical
time intervals; the saved physical time comes from the header.

The low-level reader documents native axes
\passthrough{\lstinline!(x, y, z)!} with \textbf{x vertical}. Applying
\passthrough{\lstinline!.T!} gives \passthrough{\lstinline!(z, y, x)!}:
two horizontal axes followed by height, as in the other synthesis
scripts. The line-of-sight cell length is therefore
\passthrough{\lstinline!dX[0]!}, including when the three header
spacings differ. Positive \passthrough{\lstinline!vx!} points toward an
observer at the upper x boundary, giving a blueshift.
\passthrough{\lstinline!--observer top!} implements this convention;
\passthrough{\lstinline!--observer bottom!} reverses the Doppler sign.
Only vertical, axis-aligned rays are implemented. A slanted view
requires resampling along rays and projecting all velocity components.

\passthrough{\lstinline!--downsample!} strides the two horizontal
dimensions. It does not average pixels and does not skip vertical cells
or multiply their path lengths. Rectangular horizontal domains are
supported.

The default integrates all vertical cells.
\passthrough{\lstinline!--height-start!} and
\passthrough{\lstinline!--height-stop!} select an explicit interval with
Python's inclusive-start/exclusive-stop convention. For a solar
observation, choose the lower boundary to exclude the subsurface
convection zone according to your model's geometry. Index 150 in the
existing scripts is not a universal photospheric height. There is no
automatic temperature clipping: a cool foreground structure can absorb
radiation from hotter material behind it, so dropping it changes the
question being answered. The pre-existing MURaM
\passthrough{\lstinline!tau!} field is not a SPICE line optical depth
and is not used here.

\section{Atomic populations and
abundances}\label{atomic-populations-and-abundances}

For element X and ion stage q, define

\[
A_X = n_X/n_H, \qquad F_q(T)=n_{X,q}/n_X, \qquad
p_i(T,n_e)=n_{X,q,i}/n_{X,q}.
\]

The number density of hydrogen nuclei is obtained from mass density:

\[
n_H = X_H\rho/m_H, \qquad
n_i = n_H A_X F_q(T)p_i(T,n_e).
\]

\passthrough{\lstinline!X\_H!} defaults to 0.74 and is an explicit
assumed composition parameter. Set
\passthrough{\lstinline!--hydrogen-mass-fraction!} to the value
consistent with the model EOS. The script does not infer it from an
abundance table. \passthrough{\lstinline!m\_H = 1.00794 m\_u!} is used;
atomic weights for thermal widths and all other constants are declared
in CGS in the module. No ion mass in kilograms enters a CGS
thermal-width expression.

Using rho avoids a fixed \passthrough{\lstinline!n\_H/n\_e!} conversion,
which would be particularly inappropriate in partially ionized plasma.
\passthrough{\lstinline!n\_H!} counts nuclei in all hydrogen ionization
states, not only neutral hydrogen or protons.

For each ion, the code calls
\passthrough{\lstinline!ChiantiPy.core.ion(..., temperature=T, eDensity=ne, abundance=...)!},
uses \passthrough{\lstinline!IoneqOne!} for the installed default
equilibrium balance, and calls \passthrough{\lstinline!populate()!} for
the fractional level populations. \passthrough{\lstinline!Wgfa!}
supplies transition wavelengths and Einstein A coefficients;
\passthrough{\lstinline!Elvlc!} supplies level IDs and statistical
weights \passthrough{\lstinline!mult = 2J+1!}. These interfaces are
documented in the
\href{https://github.com/chianti-atomic/ChiantiPy/blob/master/ChiantiPy/core/Ion.py}{ChiantiPy
source} and the \href{https://www.chiantidatabase.org/cug.pdf}{CHIANTI
user guide}, especially Sections 3.1 and 6.1--6.5.

The default abundance filename stem is
\passthrough{\lstinline!sun\_photospheric\_2021\_asplund!}, present in
the locally installed CHIANTI 11.0.2 database. Choose a different
installed file with \passthrough{\lstinline!--abundance!} for a
different composition or coronal FIP enhancement. Check the actual
available filenames: older distributions can differ. The selected
elemental abundance and equilibrium filename are written to metadata.
\passthrough{\lstinline!populate()!} uses its installed default
corrections and equilibrium proton density. Neither the MURaM
electron-density cube nor a CHIANTI version number alone implies that
non-equilibrium ionization or the advanced CHIANTI 11 ionization models
are being used.

\subsection{Transition identification}\label{transition-identification}

The default match window is \textbf{±0.15 Å around the paper
wavelength}. Only positive (observed-energy) wavelengths and positive A
coefficients are accepted. Theoretical wavelengths, stored as negative
numbers in CHIANTI, are excluded. If no transition matches, the run
fails rather than choosing an arbitrarily distant line. This
interpretation of the files follows
\href{https://www.chiantidatabase.org/cug.pdf}{CHIANTI's radiative-data
description}.

Every matching same-ion component is retained and its opacity/emission
is summed. This handles the closely spaced H I fine-structure components
and the two selected Fe X components at 1028.053 and 1028.082 Å in the
local data. The code prints and saves every selected wavelength, level
pair, A coefficient, statistical weight, and effective oscillator
strength. This is a feature selection window, not a full synthesis of
every line falling anywhere within the final spectral window; other-ion
blends are absent.

The local atomic-data check found C III 977.020 Å, O VI 1031.912 and
1037.613 Å, and Si VII 1049.140 Å. The latter differs from the paper by
0.110 Å and motivated the 0.15 Å matching default. In the local Si VII
file it is the
\passthrough{\lstinline!2s2 2p4 3P1!}--\passthrough{\lstinline!2s2 2p4 1S0!}
transition (levels 2--5). All 18 catalogue features were checked against
the installed database. These values are local-data findings, not
changes to the paper's wavelengths.

The nearest selected component defines the exact reference wavelength
and spectral grid center. Both this reference and the original paper
wavelength are saved. For another database release, inspect the printed
selections before increasing
\passthrough{\lstinline!--match-tolerance!}.

\subsection{Population tables}\label{population-tables}

The stored table is

\[
Q_i(\log_{10}T,\log_{10}n_e)=A_X F_q(T)p_i(T,n_e)=n_i/n_H.
\]

Only the levels needed by the selected transitions are retained. Lines
of the same ion share a solve, including the O VI doublet. Default grids
are \passthrough{\lstinline!log10(T) = 3..8!} with 101 points and
\passthrough{\lstinline!log10(ne) = 6..17!} with 45 points. The
interpolation is bilinear in the logarithms of T and ne, applied to the
linear values of Q. It therefore preserves nonnegative corner values and
handles exact zeros without taking their logarithms. Solves are skipped
only where the tabulated equilibrium ion fraction is exactly zero.

The run rejects nonfinite/negative population values, solver errors, and
populations that do not sum to unity. MURaM temperature, density, and
electron density must be positive and finite. Velocities must be finite
and below 0.01 c in magnitude. Model values outside the atomic table
cause an explicit error; there is no extrapolation, edge clipping, or
silent replacement by zero. Extend
\passthrough{\lstinline!--log-t!}/\passthrough{\lstinline!--log-ne!} to
cover your selected atmosphere. Bounds are checked while processing
column blocks.

Population caches are compressed NPZ files named by a signature. The
signature includes grid coordinates, selected transitions, ChiantiPy and
database versions, numerical abundances, the equilibrium array, and
hashes of files for the ion and next higher stage.
\passthrough{\lstinline!--rebuild-tables!} forces a rebuild. Cache
writes use an atomic rename. A cache is a performance optimization, not
an accuracy guarantee: compare maps with denser grids, particularly near
steep ion-fraction gradients. Very faint ion-fraction tails can have
large relative interpolation errors even when their absolute opacity
contribution is negligible.

\section{Opacity and spectral synthesis
equations}\label{opacity-and-spectral-synthesis-equations}

For a transition from lower l to upper u, let \(\lambda_0\) be the
physical rest wavelength in cm and \(\nu_0=c/\lambda_0\). The Doppler
width parameter is

\[
b=\sqrt{2k_BT/m_X+\xi^2}, \qquad \Delta\nu_D=\nu_0 b/c.
\]

Here \(\xi\), set by \passthrough{\lstinline!--microturbulence!}, is a
\textbf{Doppler b parameter} in km/s, not a Gaussian standard deviation
or FWHM. For this convention the Gaussian velocity standard deviation is
\(b/\sqrt{2}\). The temperature is used as both the electron temperature
for populations and the ion temperature for broadening.

With \(v_{\rm red}>0\) away from the observer, the implemented center
and profile are

\[
\nu_c=\frac{\nu_0}{1+v_{\rm red}/c}, \qquad
\phi_\nu=\frac{1}{\sqrt{\pi}\Delta\nu_D}
\exp\left[-\left(\frac{\nu-\nu_c}{\Delta\nu_D}\right)^2\right].
\]

This centers the line at \(\lambda_0(1+v_{\rm red}/c)\), consistent with
the nonrelativistic wavelength Doppler convention. Widths use the
rest-frame \(\nu_0\); relativistic transfer factors are neglected. The
profile is normalized over frequency, \(\int\phi_\nu d\nu=1\), and has
units Hz\^{}-1.

The extinction coefficient, including stimulated emission, is

\[
\boxed{\alpha_\nu=\frac{\lambda_0^2 A_{ul}}{8\pi}
\left(\frac{g_u}{g_l}n_l-n_u\right)\phi_\nu}
\quad [\mathrm{cm}^{-1}].
\]

Equivalently,

\[
\alpha_\nu=\frac{\pi e^2}{m_e c}f_{lu}
\left(n_l-\frac{g_l}{g_u}n_u\right)\phi_\nu,\qquad
f_{lu}=\frac{m_e c}{8\pi^2 e^2}\frac{g_u}{g_l}\lambda_0^2 A_{ul}.
\]

The implementation uses the A-based expression directly. The effective f
derived by this relation is also recorded; this avoids relying on a gf
field that can be absent or zero for a weak/forbidden transition. In
particular, the lower-level occupation is explicitly solved: it is not
assumed to equal the entire ion population for metastable transitions.
Negative extinction from a population inversion is retained and flagged,
rather than clipped away.

For cell-centered model data, the total column optical depth at each
frequency is

\[
\boxed{\tau(\nu)=\sum_{k=k_{\min}}^{k_{\max}-1}\alpha_{\nu,k}\,\Delta s_k},
\qquad \Delta s_k=\mathtt{dX[0]}.
\]

Summing complete cells is intentional: a trapezoidal integral over cell
centers would discard half of each boundary cell. A single-cell
atmosphere still has optical depth \(\alpha\Delta s\). The present file
format supplies uniform spacings; a genuinely nonuniform grid needs a
per-cell path-length implementation.

The spontaneous emissivity and optically thin comparison intensity are

\[
j_\nu=\frac{h\nu_0}{4\pi}A_{ul}n_u\phi_\nu,\qquad
I_{\lambda,\mathrm{thin}}=\frac{c}{\lambda^2}10^{-8}
\sum_k j_{\nu,k}\Delta s_k.
\]

The factor \(10^{-8}\) converts intensity per cm of wavelength to
intensity per Å. \textbf{Optical depth itself gets no wavelength
Jacobian}: \passthrough{\lstinline!tau2D!} simply samples \(\tau_\nu\)
at \(\nu=c/\lambda\). Both calculations include the actual column length
and number densities, with no emission-measure normalization hidden in a
contribution function.

The output does not solve \(dI_\nu/ds=-\alpha_\nu I_\nu+j_\nu\). In
particular, multiplying the total thin intensity by
\passthrough{\lstinline!exp(-tau2D)!} would incorrectly put all emission
behind a foreground screen. Emission distributed along a ray requires
depth-dependent transfer; resonance scattering additionally changes the
source function and potentially the level populations.

\section{Running the script}\label{running-the-script}

Use an environment with Python 3.9+, NumPy, SciPy, and ChiantiPy.
Matplotlib is needed for \texttt{--plot} and the standalone visualization script.
\passthrough{\lstinline!joblib!} is needed only for
\passthrough{\lstinline!--jobs!} other than 1. A CHIANTI atomic database
is a separate installation from the Python package; see the
\href{https://www.chiantidatabase.org/chianti_download.html}{official
download instructions}.

\begin{lstlisting}[language=bash]
python -m pip install numpy scipy ChiantiPy joblib matplotlib
export XUVTOP=/path/to/CHIANTI
python compute_spice_optical_depth.py --list-lines
\end{lstlisting}

The repository's \passthrough{\lstinline!.venv!} had none of these
scientific packages during validation. An existing environment with them
was available at
\passthrough{\lstinline!/Users/memolnar/miniconda3/envs/SPICE\_DEMs/bin/python!},
with the database at
\passthrough{\lstinline!/Users/memolnar/Work/CHIANTI!}. No environments
or database files were modified.

Start with a single C III snapshot and a sparse horizontal sample. The
paths and iteration below follow
\passthrough{\lstinline!example\_CIII.py!}; adapt them to your machine:

\begin{lstlisting}[language=bash]
python compute_spice_optical_depth.py \
  --model-dir /glade/campaign/hao/radmhd/Rempel/Spot_Motion/AR_25Mm_v1000_vz300_disp/3D \
  --iterations 400000 \
  --lines c_3_977 \
  --height-start 150 \
  --downsample 8 \
  --output-dir optical_depth_CIII \
  --jobs 8
\end{lstlisting}

Here height index 150 is an example borrowed from the repo's existing
boundary choice; establish that it is suitable for this snapshot before
interpreting the column. \passthrough{\lstinline!--height-stop!}
defaults to the top of the model. Leaving out
\passthrough{\lstinline!--height-start!} includes the entire input
column.

Run all eight default lines over an iteration range:

\begin{lstlisting}[language=bash]
python compute_spice_optical_depth.py \
  --model-dir /path/to/MURaM/3D \
  --t-start 400000 --t-end 502000 --time-step 2000 \
  --height-start 150 --downsample 2 \
  --output-dir optical_depth_default --jobs 8
\end{lstlisting}

Or explicitly choose the resonance diagnostics:

\begin{lstlisting}[language=bash]
python compute_spice_optical_depth.py \
  --model-dir /path/to/MURaM/3D --iterations 400000 402000 \
  --lines c_3_977 o_6_1031 o_6_1037 ne_8_770 \
  --height-start 150 --microturbulence 10 \
  --velocity-range 250 --num-waves 1001
\end{lstlisting}

To compute the whole catalogue, substitute
\passthrough{\lstinline!--all-lines!} for
\passthrough{\lstinline!--lines ...!}. For a cooler/denser selected
atmosphere or a convergence check, change the table grids, for example
\passthrough{\lstinline!--log-t 3 8 201 --log-ne 5 18 105!}.

Build tables on their own, reusable in subsequent matching runs:

\begin{lstlisting}[language=bash]
python compute_spice_optical_depth.py --tables-only \
  --lines c_3_977 o_6_1031 o_6_1037 ne_8_770 \
  --cache-dir opacity_tables --jobs 8
\end{lstlisting}

Table creation can be expensive for ions with many levels. Processes
parallelize the density samples; snapshots and line features are
processed sequentially. On HPC systems, constrain BLAS threads (for
example \passthrough{\lstinline!OPENBLAS\_NUM\_THREADS=1!} and
\passthrough{\lstinline!OMP\_NUM\_THREADS=1!}) when using many worker
processes, and allocate enough memory for each worker.
\passthrough{\lstinline!--jobs -1!} requests all available CPUs.

Input cubes stay memory mapped.
\passthrough{\lstinline!--chunk-columns!} limits the atmospheric arrays
held for integration, and wavelengths are streamed without allocating a
full columns × heights × wavelengths profile array. Output arrays remain
in memory for one line and snapshot: the two large float64 arrays
together need about
\passthrough{\lstinline!16 * n\_horizontal\_1 * n\_horizontal\_2 * num\_waves!}
bytes, plus maps, table and compression overhead. A 512 × 512 field with
401 wavelengths needs about 1.57 GiB for those two arrays alone. Reduce
horizontal sampling or wavelength count when necessary; chunking alone
does not reduce output-array memory.

Existing output files are refused unless
\passthrough{\lstinline!--overwrite!} is supplied. Different abundances,
height cuts, or broadening choices should normally use different output
directories so they can be compared. Files required for every selected
iteration are checked before building tables.

\section{Outputs and interpretation}\label{outputs-and-interpretation}

One file is saved per feature and snapshot:
\passthrough{\lstinline!spice\_opacity\_c\_3\_977\_400000.npz!}, for
example.

{\def\LTcaptype{none} % do not increment counter
\begin{longtable}[]{@{}
  >{\raggedright\arraybackslash}p{(\linewidth - 4\tabcolsep) * \real{0.3333}}
  >{\raggedright\arraybackslash}p{(\linewidth - 4\tabcolsep) * \real{0.3333}}
  >{\raggedright\arraybackslash}p{(\linewidth - 4\tabcolsep) * \real{0.3333}}@{}}
\toprule\noalign{}
\begin{minipage}[b]{\linewidth}\raggedright
Key
\end{minipage} & \begin{minipage}[b]{\linewidth}\raggedright
Shape
\end{minipage} & \begin{minipage}[b]{\linewidth}\raggedright
Meaning
\end{minipage} \\
\midrule\noalign{}
\endhead
\bottomrule\noalign{}
\endlastfoot
\passthrough{\lstinline!tau2D!} &
\texttt{(1, native\_z, native\_y, }\allowbreak\texttt{n\_wavelengths)} &
Total selected-column spectral optical depth \\
\passthrough{\lstinline!I2D!} & same & Optically thin intensity, erg
s\^{}-1 cm\^{}-2 sr\^{}-1 Å\^{}-1 \\
\passthrough{\lstinline!tau\_rest!} &
\passthrough{\lstinline!(1, native\_z, native\_y)!} & Optical depth at
the exact atomic reference wavelength, including model velocities \\
\passthrough{\lstinline!tau\_peak!} & same map shape & Maximum optical
depth among the sampled wavelengths \\
\passthrough{\lstinline!tau\_static!} & same map shape & Optical depth
at the reference wavelength with all resolved velocities set to zero \\
\passthrough{\lstinline!wvlGrid!} &
\passthrough{\lstinline!(n\_wavelengths,)!} & Increasing physical
wavelengths in Å \\
\passthrough{\lstinline!velocityGrid!} &
\passthrough{\lstinline!(n\_wavelengths,)!} & Offsets relative to the
atomic reference in km/s, red positive \\
\passthrough{\lstinline!detectorWvlGrid!} &
\passthrough{\lstinline!(n\_wavelengths,)!} &
\passthrough{\lstinline!order * wvlGrid!}, a coordinate only, not an
instrument response \\
\passthrough{\lstinline!metadata!} & scalar JSON string & Atomic
selections, units, parameters, versions, file hashes, geometry,
diagnostics and command \\
\end{longtable}
}

The leading dimension is length one for compatibility with the
repository's time-leading \passthrough{\lstinline!I2D!} convention.
Horizontal dimensions include any downsampling. These files can supply
\passthrough{\lstinline!I2D!} and \passthrough{\lstinline!wvlGrid!} to
the existing visualization script, but their intensity model must be
kept in mind. The metadata loads with
\passthrough{\lstinline!allow\_pickle=False!}; no Python object
deserialization is required.

\passthrough{\lstinline!tau\_static!} is a useful comparison when
velocity gradients move absorbing layers out of resonance. It includes
thermal and selected microturbulent broadening. For a multiplet it is
the sum of component opacities at the \textbf{shared reference
wavelength}, not a sum of each component's individual peak.

\passthrough{\lstinline!tau\_peak!} depends on the chosen finite
spectral grid and is not a guaranteed continuous maximum. The script
warns when resolved velocities exceed the spectral window or when the
sample spacing is greater than one third of the smallest Doppler b.
These are basic checks: extend the window enough to capture the wings
and repeat with finer sampling until the maps and integrated intensities
converge. Even in the absence of a warning, a near-edge line can lose
wing flux. The odd default of 401 samples guarantees the exact reference
wavelength is present.

\subsection{Matplotlib maps of maximum column optical depth}

Add \texttt{--plot} to a synthesis command to save a PNG beside every NPZ:

\begin{lstlisting}[language=bash]
python compute_spice_optical_depth.py \
  --model-dir /path/to/MURaM/3D --iterations 400000 \
  --lines c_3_977 --height-start 150 --downsample 8 --plot
\end{lstlisting}

The image is named \texttt{spice\_opacity\_c\_3\_977\_400000\_tau\_peak.png}.
To visualize already saved results without repeating synthesis:

\begin{lstlisting}[language=bash]
python visualize_spice_optical_depth.py optical_depth_output/spice_opacity_*.npz \
  --output-dir optical_depth_maps
\end{lstlisting}

The standalone command reads only \texttt{tau\_peak} and \texttt{metadata},
and requires NumPy and Matplotlib. It does not load the spectral cubes,
CHIANTI, or MURaM inputs. Use \texttt{--overwrite} to replace existing images.
Both commands support \texttt{--plot-scale log|linear},
\texttt{--plot-limits MIN MAX}, \texttt{--plot-cmap}, and \texttt{--plot-dpi}.
Defaults are logarithmic grayscale, fixed optical-depth limits of \texttt{1e-3}
to \texttt{5}, and 180 DPI. Values at or below \texttt{1e-3} are black and
values at or above \texttt{5} are white; the colorbar spans those same limits
and has the map's height. Limits are optical depths, not their logarithms.
Use \texttt{--plot-limits} or \texttt{--plot-cmap} to choose a different
display scale.

Each pixel shows
\[
\tau_{\mathrm{peak}}(x,y)=\max_\lambda\int_{\mathrm{LOS}}\alpha_\lambda\,ds.
\]
For nonnegative extinction, cumulative optical depth is largest through the
complete selected column. The map then takes its maximum over the sampled
wavelengths; it is not a maximum of the local extinction coefficient. It does
not calculate a depthwise maximum of cumulative optical depth for inverted
populations, where extinction can be negative.

The plot labels the two horizontal coordinates \textbf{x [Mm]} and
\textbf{y [Mm]}: plot x is native MURaM y, plot y is native MURaM z,
and native MURaM x remains the line of sight. Stored maps have row/column
order \texttt{(native\_z, native\_y)}, so they are plotted without transposition.
Coordinates are relative to the domain's horizontal lower boundaries.
Sampled cell centers are calculated as
\texttt{(native\_index + 0.5) * cell\_spacing\_cm / 1e8}; display edges lie
halfway between sampled columns and end at the domain boundaries. Unequal
spacings, rectangular domains, and downsampling factors that do not divide
the domain size are supported. Colors represent nearest sampled columns,
not pixel averages.

A white contour marks \(\tau_{\mathrm{peak}}=1\) when the finite map spans
that value and both horizontal dimensions have multiple samples.
In logarithmic maps, finite values at or below the lower limit, including zero
and negative values, are black; nonfinite pixels are gray and their count is
annotated. Linear colors retain finite negative optical depths. Rendering uses
a noninteractive Matplotlib canvas, so synthesis on an HPC compute node does
not open windows.

Read a result and use the plotting helper from Python:

\begin{lstlisting}[language=Python]
import json
import numpy as np
from scipy.integrate import trapezoid
from visualize_spice_optical_depth import save_tau_peak_map

with np.load("optical_depth_CIII/spice_opacity_c_3_977_400000.npz",
             allow_pickle=False) as data:
    tau_peak = data["tau_peak"][0]
    tau_static = data["tau_static"][0]
    integrated_thin = trapezoid(data["I2D"][0], data["wvlGrid"], axis=-1)
    metadata = json.loads(str(data["metadata"]))

print("Fraction of sampled columns with tau_peak >= 1:", (tau_peak >= 1).mean())
print("Atomic components:", metadata["atomic"]["components"])
save_tau_peak_map(tau_peak, metadata, "CIII_tau_peak.png", limits=(0.01, 100))
\end{lstlisting}

For non-inverted populations, \passthrough{\lstinline!tau << 1!} means
weak line extinction along the selected column, while
\passthrough{\lstinline!tau >= 1!} indicates potentially consequential
opacity. The total optical depth through a column can exceed the optical
depth between a particular emitting layer and the observer. These
products do not locate a tau-one surface or save cumulative
depth-dependent optical depths.

The O VI doublet provides a useful internal comparison. With negligible
upper populations, equal lower-level columns, and the same velocity
broadening, its static peak-depth ratio scales as
\passthrough{\lstinline!f * wavelength!}, approximately two. It is not
an assumption enforced by the code, and observed intensity ratios also
depend on excitation and transfer.

\section{Limits and validation}\label{limits-and-validation}

The calculation assumes a Maxwellian electron distribution, equilibrium
ion fractions, statistical-equilibrium level populations, a single
ion/electron temperature, and Gaussian profiles. It includes neither an
external pumping radiation field nor a self-consistent scattering source
function. Natural and pressure wings, partial redistribution,
non-equilibrium ionization, and instrumental PSFs/effective areas are
absent. It is not a calculation of SPICE detector counts. Absence of
line opacity does not by itself establish that an observed spectrum is
described fully by collisional emission.

Continuum opacity is also absent. In particular, neutral-hydrogen
bound-free absorption can affect the short-wavelength band below the
Lyman edge near 912 Å and the physical Si XII wavelength. A full
foreground absorption model would need the appropriate neutral
populations and continuum cross sections. The saved optical depths refer
exclusively to the selected bound-bound components.

The test suite uses analytic slabs and artificial MURaM-format files. It
checks absolute line-center normalization, total cell length (including
a one-cell column), stimulated emission and inversions, thermal and
microturbulent widths, Doppler shifts, multiplet addition, profile
normalization, and the wavelength emissivity Jacobian. A nonsquare
fixture with unequal grid spacings checks the reader transpose, vertical
path length, height slicing, and observer sign. Invalid cells,
out-of-grid values, missing inputs, and output collisions are also
exercised.

\begin{lstlisting}[language=bash]
python -m unittest discover -s tests -v
SINSPICE_TEST_CHIANTI=1 python -m unittest discover -s tests -v
\end{lstlisting}

The second command additionally checks actual catalogue identifications
and builds a small C III table, compares a tabulated population against
a direct ChiantiPy solve, verifies cache reuse, and runs the CLI
workflow on a generated MURaM-format snapshot, including its PNG output.
Plotting tests check the map orientation and physical extent with unequal
spacings and non-divisible downsampling, zero/negative/nonfinite pixels, and
standalone plotting from an NPZ containing only the map and metadata.
All 15 tests passed using
ChiantiPy 0.15.2 and CHIANTI 11.0.2. Separate smoke runs built tables
for all 18 features on a 2 × 2 grid and built the full default 101 × 45
C III table with two worker processes. The tiny smoke-test grids are
deliberately insufficient for science; they exercise the interfaces and
solver paths. Artificial fixtures are not solar-model validation.
Production MURaM snapshots at the hard-coded GLADE paths were not
available locally, so no physical opacity maps for those simulations are
claimed here.

For a scientific production run, first establish the correct geometry
and composition, then converge the population grids and spectral
sampling. Use non-LTE/NEI calculations where the inferred optical depths
or timescales demand them. Record the saved atomic metadata alongside
any reported results.

\end{document}
